% !TEX program = pdflatex
% ======================================================================
%  PFA : compte rendu d'avancement, fin de la semaine 2 (livrable L2)
%  Module : Développement et déploiement d'applications intelligentes
%
%  MODE D'EMPLOI
%  1. Renseignez le bloc « Informations de l'équipe » ci-dessous.
%  2. Remplacez chaque zone \zone{...} par votre réponse.
%  3. Cochez une case en remplaçant \case par \coche.
%  4. Passez \consignestrue à \consignesfalse avant de rendre le document.
%  5. Compilez avec pdflatex (deux passes) et déposez le PDF dans le dépôt :
%     docs/compte-rendu-semaine2.pdf
% ======================================================================
\documentclass[11pt,a4paper]{article}

\usepackage[utf8]{inputenc}
\usepackage[T1]{fontenc}
\usepackage{lmodern}
% Le français est chargé s'il est installé (Overleaf, TeX Live complet).
\IfFileExists{french.ldf}{%
  \usepackage[french]{babel}\frenchsetup{StandardLists=true}}{}
\usepackage[margin=2.2cm,headheight=15pt]{geometry}
\usepackage{xcolor}
\usepackage{array,tabularx,booktabs}
\usepackage{enumitem}
\usepackage{fancyhdr}
\usepackage{amssymb}
\usepackage{titlesec}
\usepackage{needspace}
\usepackage[most]{tcolorbox}
\usepackage[hidelinks]{hyperref}

% ---------------------------------------------------------------------
%  Informations de l'équipe (À REMPLIR)
% ---------------------------------------------------------------------
\newcommand{\NomCourt}{Assistant RAG conformité}
\newcommand{\NumProjet}{n\textsuperscript{o}~19}
\newcommand{\IntituleProjet}{Assistant documentaire RAG pour la conformité des documents publics marocains}
\newcommand{\Depot}{https://github.com/joshBM7/assistant-documentaire-rag-conformite-marocaine}
\newcommand{\CommitRef}{\texttt{2e8b6c6} (livrable L2)}
\newcommand{\DateRemise}{27/11/2026}
\newcommand{\Coordinateur}{Joshua-Phil-Lilan BAYONNE-MANGA}

% Affichage des consignes : \consignesfalse avant de rendre le document
\newif\ifconsignes
\consignesfalse

% ---------------------------------------------------------------------
%  Mise en forme
% ---------------------------------------------------------------------
\definecolor{encre}{HTML}{12343B}
\definecolor{accent}{HTML}{0F766E}
\definecolor{fond}{HTML}{EAF3F2}
\definecolor{aremplir}{HTML}{B45309}

\newcommand{\zone}[1]{\textcolor{aremplir}{\itshape #1}}
\newcommand{\case}{$\square$}
\newcommand{\coche}{$\boxtimes$}

\newcolumntype{L}{>{\raggedright\arraybackslash}X}
\newcolumntype{P}[1]{>{\raggedright\arraybackslash}p{#1}}
\renewcommand{\arraystretch}{1.35}
\setlength{\parindent}{0pt}
\setlength{\parskip}{6pt}
\setlist{nosep,leftmargin=1.4em}

\titleformat{\section}{\large\bfseries\color{encre}}{\thesection.}{0.5em}{}
\titlespacing*{\section}{0pt}{14pt}{4pt}

\pagestyle{fancy}
\fancyhf{}
\fancyhead[L]{\small\color{encre} PFA, compte rendu d'avancement, semaine 2}
\fancyhead[R]{\small\color{encre}\NomCourt}
\fancyfoot[C]{\small\thepage}
\renewcommand{\headrulewidth}{0.4pt}

\newtcolorbox{consigne}{colback=fond,colframe=fond,boxrule=0pt,arc=2mm,
  left=3mm,right=3mm,top=2mm,bottom=2mm,fontupper=\small}

\hypersetup{pdftitle={PFA, compte rendu d'avancement, semaine 2},pdfauthor={\NomCourt}}

% ======================================================================
\begin{document}

\begin{center}
  {\small\color{accent}\bfseries Développement et déploiement d'applications intelligentes}\\[4pt]
  {\LARGE\bfseries\color{encre} Compte rendu d'avancement du PFA}\\[4pt]
  {\large\color{encre} Fin de la semaine 2 : livrable L2, « MVP mesuré »}
\end{center}

\ifconsignes
\begin{consigne}
\textbf{Consignes (à masquer avant la remise avec \texttt{\textbackslash consignesfalse}).}
Ce compte rendu justifie l'avancement du projet à la fin de la semaine 2. Il sert de support
à la note de suivi de cette semaine (10 points) : livrable démontré (5 points), compte rendu
tenu et honnête (2,5 points), contribution par membre (2,5 points).
\begin{itemize}
  \item \textbf{Longueur :} 4 pages maximum, hors annexes. Des phrases courtes, des chiffres, des liens.
  \item \textbf{Preuves :} chaque affirmation renvoie à un commit, une \emph{pull request}, une
        exécution de la CI ou une capture. Une ligne sans preuve compte comme non faite.
  \item \textbf{Honnêteté :} un retard expliqué vaut mieux qu'un retard caché. Les écarts se
        déclarent à la rubrique~\ref{sec:ecarts}.
  \item \textbf{Remise :} PDF déposé dans le dépôt, fichier \texttt{docs/compte-rendu-semaine2.pdf},
        au plus tard à la fin de la semaine 2.
\end{itemize}
\end{consigne}
\fi

% ---------------------------------------------------------------------
\section{Identification}

\begin{tabularx}{\linewidth}{@{}P{4.6cm}L@{}}
\toprule
Projet & \NumProjet{} : \IntituleProjet \\
Nom court & \NomCourt \\
Dépôt GitHub & \url{\Depot} \\
Version de référence du livrable & \CommitRef \\
Coordinateur & \Coordinateur \\
Date de remise & \DateRemise \\
\bottomrule
\end{tabularx}

\medskip
\begin{tabularx}{\linewidth}{@{}LP{4.2cm}P{4.2cm}@{}}
\toprule
\textbf{Membre (nom et prénom)} & \textbf{Compte GitHub} & \textbf{Rôle principal} \\
\midrule
Joshua-Phil-Lilan BAYONNE-MANGA & \texttt{joshBM7} & Coordination, architecture RAG, intégration \& livrable \\
Djamila HACHIMI SOBA & Compte GitHub à confirmer & Corpus, données, validation métier \\
\bottomrule
\end{tabularx}

% ---------------------------------------------------------------------
\section{Synthèse de l'avancement}

\textbf{État global à la fin de la semaine 2 :}\\
\mbox{\case~en avance} \quad \mbox{\coche~conforme au planning} \quad \mbox{\case~en retard rattrapable} \quad \mbox{\case~en retard critique}

\textbf{En cinq lignes au plus :} ce qui fonctionne aujourd'hui, ce qui manque, et le principal
risque pour la semaine 3.

Le corpus de 7 textes marocains est ingéré et 649 passages sont indexés. La recherche hybride,
les garde-fous de citation, le contrôle de conformité et les endpoints FastAPI sont implémentés.
Les 36 tests automatisés passent sans téléchargement ni appel externe. Il reste à valider le parcours
réel avec la clé LLM, à produire au moins 20 cas d'évaluation et à mesurer les métriques.
Risque principal : la qualité et la latence avec les vrais modèles ne sont pas encore mesurées.

% ---------------------------------------------------------------------
\needspace{7cm}
\section{Rappel du périmètre validé}

\textbf{Scénario principal} (une entrée, un résultat visible) : poser une question réglementaire
sur un domaine précis; l'agent recherche dans le corpus, appelle le modèle, cite les passages
avec \texttt{[S\#]} et retourne une réponse vérifiable.

\textbf{Scénario dégradé} (abstention, demande de précision ou passage à un humain) : si aucune
source pertinente n'est trouvée, si une citation est absente ou si un article est inventé, le
système refuse la réponse et explique la cause.

\begin{tabularx}{\linewidth}{@{}P{0.7cm}LP{2.6cm}P{2.3cm}P{3.6cm}@{}}
\toprule
\textbf{N\textsuperscript{o}} & \textbf{Fonctionnalité du MVP} & \textbf{Responsable} & \textbf{État} & \textbf{Preuve} \\
\midrule
1 & Ingestion de 7 textes et découpage par article & Équipe & Fait & \texttt{ragconf/ingest.py}, \texttt{ragconf/chunking.py}, 649 passages indexés \\
2 & Recherche hybride dense + BM25 et filtres & Équipe & Fait & \texttt{ragconf/retrieve.py}, tests de récupération \\
3 & Réponse citée avec garde-fous & Équipe & Fait & \texttt{ragconf/pipeline.py}, \texttt{ragconf/guardrails.py} \\
4 & Vérification de conformité par exigences & Équipe & Partiel & \texttt{ragconf/compliance.py}; jeu de cas réel à finaliser \\
5 & API FastAPI et interface Streamlit & Équipe & Partiel & \texttt{ragconf/api.py}, \texttt{app.py}; démonstration avec clés externes à faire \\
\bottomrule
\end{tabularx}

\textbf{Le périmètre a-t-il changé depuis la fiche de cadrage ?}\\
\mbox{\coche~non} \quad \mbox{\case~oui, changement approuvé le 27/11/2026} \quad
\mbox{\case~oui, demande à valider (rubrique~\ref{sec:plan})}

% ---------------------------------------------------------------------
\section{Contenu du livrable L2}

Cochez l'état de chaque élément et donnez la preuve. Un élément « partiel » précise ce qui manque.

\begin{tabularx}{\linewidth}{@{}LP{0.9cm}P{1.4cm}P{1.3cm}P{4.2cm}@{}}
\toprule
\textbf{Élément attendu} & \textbf{Fait} & \textbf{Partiel} & \textbf{Non fait} & \textbf{Preuve (lien, commande, capture)} \\
\midrule
Scénario principal complet, sans simulation sur le chemin principal & \coche & \case & \case & Code complet; démonstration réelle avec clé LLM encore à exécuter \\
Au moins deux outils réels appelés par l'agent & \case & \coche & \case & Recherche BM25/dense et extraction-conformité implémentées; qualification « agent » à démontrer \\
Garde-fous actifs : abstention ou escalade, action sensible simulée ou validée par un humain & \coche & \case & \case & \texttt{ragconf/guardrails.py}; refus testé par \texttt{pytest} \\
API FastAPI : \texttt{/health} et un \emph{endpoint} métier & \coche & \case & \case & \texttt{/health}, \texttt{/documents}, \texttt{/ask}, \texttt{/compliance}; tests TestClient \\
Tests \texttt{pytest} qui passent & \coche & \case & \case & 36 tests passent (`python -m pytest -q`) \\
CI verte sur la branche principale & \case & \case & \coche & Aucun workflow CI configuré dans \texttt{.github/workflows/} \\
\texttt{Dockerfile} qui se construit & \case & \case & \coche & Construction lancée mais Docker Desktop indisponible; fichiers présents \\
Jeu de test porté à 20 cas & \case & \case & \coche & Modèle actuel : 4 cas seulement dans \texttt{data/eval/questions.template.jsonl} \\
Premier score mesuré sur le jeu de test & \case & \case & \coche & Aucun jeu rempli ni métrique calculée \\
Trois échecs analysés & \coche & \case & \case & Cas limites et corrections dans \texttt{rubrique~\ref{sec:echecs}} \\
\bottomrule
\end{tabularx}

% ---------------------------------------------------------------------
\section{Démonstration reproductible}

\textbf{Commandes pour reproduire la démonstration} à partir d'un clone du dépôt :

\begin{tcolorbox}[colback=white,colframe=encre!40,boxrule=0.4pt,arc=1mm,
  left=2mm,right=2mm,top=1mm,bottom=1mm,fontupper=\small\ttfamily]
git clone \Depot\\
cd assistant-documentaire-rag-conformite-marocaine\\
python -m venv .venv \&\& .venv\textbackslash Scripts\textbackslash activate\\
pip install -r requirements.txt \&\& cp .env.example .env\\
python -m ragconf.cli ingest\\
uvicorn ragconf.api:app --port 8000
\end{tcolorbox}

\needspace{3.5cm}
\textbf{Un exemple réel} tiré du jeu de test :

\begin{tabularx}{\linewidth}{@{}P{3.2cm}L@{}}
\toprule
Entrée & ``python -m ragconf.cli ingest'' \\
Résultat obtenu & 649 passages indexés à partir de 7 textes PDF \\
Outils appelés & Extraction PyMuPDF/OCR, nettoyage, découpage par article, embeddings/ChromaDB \\
Durée de la réponse & À mesurer sur une machine vierge \\
\bottomrule
\end{tabularx}

% ---------------------------------------------------------------------
\section{Résultats mesurés}\label{sec:resultats}

Reprenez les critères de réussite de la fiche de cadrage (rubrique F). Les valeurs sont mesurées,
pas estimées.

\begin{tabularx}{\linewidth}{@{}LP{2.2cm}P{2.6cm}P{4.4cm}@{}}
\toprule
\textbf{Indicateur} & \textbf{Cible} & \textbf{Valeur mesurée} & \textbf{Méthode de mesure} \\
\midrule
Passages indexés / corpus complet & 649 passages & 649 passages & `python -m ragconf.cli ingest`; 7 PDF dans `data/raw/` \\
Tests automatisés réussis & 36 sur 36 & 36 sur 36 & `python -m pytest -q` \\
Taux de réponses correctes & 80 \% & Non mesuré & Jeu de 20 cas à construire dans `data/eval/` \\
\bottomrule
\end{tabularx}

\textbf{Jeu de test :} 4 modèles de cas seulement, dont 1 hors périmètre; aucun jeu de 20 cas rempli.
\textbf{Cas réussis :} non mesuré; les 36 tests unitaires du dépôt passent.

% ---------------------------------------------------------------------
\section{Analyse de trois échecs}\label{sec:echecs}

\begin{tabularx}{\linewidth}{@{}P{2.6cm}P{2.6cm}P{2.6cm}LL@{}}
\toprule
\textbf{Cas (entrée)} & \textbf{Attendu} & \textbf{Obtenu} & \textbf{Cause probable} & \textbf{Correction prévue} \\
\midrule
Réponse avec article 99 non présent & Refus ou citation vérifiable & Article inventé potentiellement cité & Hallucination ou validation trop faible & Garde-fou vérifiant les citations et articles avant réponse \\
Question hors corpus & Réponse ``information non disponible'' & Proposition speculative & Aucun contexte assez pertinent & Refus par défaut et demande de précision \\
Exigence reformulée dans le document & Verdict ``non déterminé'' & Rapport conforme ou non conforme à tort & Recherche BM25 trop littérale & Combiner plusieurs méthodes de recherche et revue humaine \\
\bottomrule
\end{tabularx}

% ---------------------------------------------------------------------
\section{Prévu, réalisé, écarts}\label{sec:ecarts}

\begin{tabularx}{\linewidth}{@{}LLL@{}}
\toprule
\textbf{Prévu pour la semaine 2} & \textbf{Réalisé} & \textbf{Écart et cause} \\
\midrule
Compléter la fiche de cadrage & Fiche remplie dans \texttt{Fiche\_de\_cadrage\_completee.md} & Aucun écart majeur \\
Ingestion du corpus & 7 PDF lus et 649 passages indexés & Données présentes localement; URLs/licences encore à compléter \\
API, interface, garde-fous & Implémentés et testés avec faux modèles & Démonstration avec clé externe et modèles réels non réalisée \\
Évaluation à 20 cas & Modèle à 4 cas seulement & Écart prévu : rédaction des 20 cas reportée à la semaine 3 \\
\bottomrule
\end{tabularx}

\textbf{Difficultés rencontrées et parades :}

\begin{tabularx}{\linewidth}{@{}LP{2.4cm}L@{}}
\toprule
\textbf{Difficulté ou risque} & \textbf{Gravité} & \textbf{Parade appliquée ou prévue} \\
\midrule
Accès à une clé LLM réelle et limites de coût & forte & Tests hors ligne conservés; validation réelle planifiée avec quota limité \\
Docker Desktop indisponible sur la machine & moyenne & Vérification Docker à reprendre sur machine disponible; commandes documentées \\
\bottomrule
\end{tabularx}

% ---------------------------------------------------------------------
\section{Contribution de chaque membre}

Le contenu de ce tableau doit correspondre à l'historique du dépôt.

\begin{tabularx}{\linewidth}{@{}P{3cm}LP{3.4cm}P{1.8cm}@{}}
\toprule
\textbf{Membre} & \textbf{Travail réalisé cette semaine} & \textbf{Commits et PR (liens)} & \textbf{Part estimée} \\
\midrule
Joshua-Phil-Lilan BAYONNE-MANGA & Architecture, pipeline RAG, API, interface, ingestion, garde-fous et fiche de suivi & Dépôt, commit \texttt{17b4bb5} et historique Git & 60 \% \\
Djamila HACHIMI SOBA & Corpus de textes, inventaire des données, validation métier et cadrage & À compléter avec les commits/PR de la membre & 40 \% \\
\bottomrule
\end{tabularx}

\textbf{Fonctionnement de l'équipe cette semaine} (réunions tenues, revue des \emph{pull requests},
désaccords éventuels) : Répartition documentée dans l'historique Git; revue croisée des modifications
avant la remise prévue pour la semaine 3.

% ---------------------------------------------------------------------
\section{Plan de la semaine 3}\label{sec:plan}

Rappel : la semaine 3 sert à mesurer, corriger, documenter et préparer la soutenance, pas à ajouter
des fonctionnalités. Le code est gelé 48 heures avant la soutenance.

\begin{tabularx}{\linewidth}{@{}LP{2.6cm}P{2.2cm}L@{}}
\toprule
\textbf{Tâche} & \textbf{Responsable} & \textbf{Échéance} & \textbf{Critère de « terminé »} \\
\midrule
Construire 20 cas d'évaluation réels & Équipe & 30/11/2026 & 20 cas JSONL remplis, catégories et réponses attendues présentes \\
Mesurer les métriques cibles & Équipe & 01/12/2026 & Rapport chiffré : exactitude, rappel, citations et refus \\
Démontrer le parcours RAG réel & Joshua-Phil-Lilan BAYONNE-MANGA & 02/12/2026 & Une question, une réponse citée et une abstention reproductibles \\
Valider Docker et préparer la soutenance & Équipe & 03/12/2026 & Docker Compose lancé sur une machine disponible; démonstration gelée \\
\bottomrule
\end{tabularx}

\textbf{Modification de périmètre demandée :} aucune; ne pas étendre la couverture réglementaire pendant la semaine 3.

\textbf{Question ou demande à l'enseignant :} confirmation du niveau d'exigence pour la mention « agent » : routeur explicite entre question, recherche et refus est-il suffisant ?

% ---------------------------------------------------------------------
\section{Déclaration de l'équipe}

Nous certifions que ce compte rendu décrit l'état réel du projet à la date de remise et que la
répartition du travail indiquée est exacte.

\begin{tabularx}{\linewidth}{@{}LP{3cm}P{5cm}@{}}
\toprule
\textbf{Membre} & \textbf{Date} & \textbf{Signature} \\
\midrule
Joshua-Phil-Lilan BAYONNE-MANGA & 27/11/2026 & \\[10pt]
Djamila HACHIMI SOBA & & \\[10pt] \\[10pt]
\bottomrule
\end{tabularx}

% ---------------------------------------------------------------------
\appendix
\section*{Annexes (hors limite de pages)}

Insérez ici, avec une légende pour chacune : la capture de la CI verte, une capture de la
démonstration, le tableau complet des résultats sur le jeu de test.

% Exemple d'insertion d'une capture (ajoutez \usepackage{graphicx} dans le préambule) :
% \begin{center}
%   \includegraphics[width=0.9\linewidth]{figures/ci-verte.png}\\
%   {\small Figure A1 : exécution de la CI sur la branche principale.}
% \end{center}

\end{document}
